%HW08.tex
%
% Eighth Homework for Honors Algebra I  (415H)
% Frank Sottile
%%%%%%%%%%%%%%%%%%%%%%%%%%%%%%%%%%%%%%%%%%%%%%%%%%%%%%%%%%%%%%%%%%%%%%%
\documentclass[12pt]{article}
\usepackage{multicol,amsfonts, amssymb,  mathtools,amsmath}
\usepackage[dvipsnames]{xcolor}
\usepackage{colordvi,graphicx}
\headheight=8pt
%
\topmargin=-75pt
\textheight=720pt   \textwidth=560pt
\oddsidemargin=-60pt \evensidemargin=-60pt

\pagestyle{empty}   %  Change this for your write-up

%%%%%%%%%%%%%%%%%%%%%%%%%%%%%%%%%%%%%%%%%%%%

\newcommand{\CC}{{\mathbb C}}
\newcommand{\NN}{{\mathbb N}}
\newcommand{\QQ}{{\mathbb Q}}
\newcommand{\RR}{{\mathbb R}}
\newcommand{\ZZ}{{\mathbb Z}}

\newcommand{\Aut}{\textrm{Aut}}
\newcommand{\op}{\textrm{op}}

\newcommand{\vect}[2]{(\begin{smallmatrix}#1\\#2\end{smallmatrix})}

\def\Color#1#2{\special{color push cmyk #1}#2\special{color pop}}
%\def\Indigo#1{\Color{.42 1. 0. .49}{#1}}
\def\Indigo#1{\Color{1. .95 .05 .4}{#1}}
\def\MyViolet#1{\Color{.6 1. 0. .15}{#1}}
\def\TAMU#1{\Color{.15 1. .39 .49}{#1}}

\newcommand{\barsl}{\noindent\begin{minipage}[t]{590pt}
\Indigo{\rule{590pt}{1.2pt}}\vspace{-5.7mm}\\
\MyViolet{\rule{590pt}{1.2pt}}\vspace{-5.7mm}\\
\Blue{\rule{590pt}{1.2pt}}\vspace{-5.7mm}\\
\Green{\rule{590pt}{1.2pt}}\vspace{-5.7mm}\\
\Yellow{\rule{590pt}{1.2pt}}\vspace{-5.7mm}\\
\Orange{\rule{590pt}{1.2pt}}\vspace{-5.7mm}\\
\Red{\rule{590pt}{1.2pt}}\bigskip
\end{minipage}}


\newcommand{\barsn}{\noindent\begin{minipage}[t]{560pt}
\Indigo{\rule{560pt}{1.1pt}}\vspace{-4.5mm}\\
\MyViolet{\rule{560pt}{1.1pt}}\vspace{-4.5mm}\\
\Blue{\rule{560pt}{1.1pt}}\vspace{-4.5mm}\\
\Green{\rule{560pt}{1.1pt}}\vspace{-4.5mm}\\
\Yellow{\rule{560pt}{1.1pt}}\vspace{-4.5mm}\\
\Orange{\rule{560pt}{1.1pt}}\vspace{-4.5mm}\\
\Red{\rule{560pt}{1.1pt}}\bigskip
\end{minipage}}

\def\demph#1{\TAMU{{\sl #1}}}
\def\defcolor#1{\TAMU{#1}}

\begin{document}
\LARGE 
\noindent
Algebra 415H \ \ Autumn 2026\vspace{2pt}\\
{\color{magenta}Your Name}\vspace{2pt}\\
\Large 15 October 2026 \hfill
\sf
 Eighth Homework\makebox[40pt][l]{\ }
\normalsize\vspace{10pt}

\noindent
Write your answers neatly, in complete sentences.  
I highly recommend recopying your work before handing it in.
Correct and crisp proofs are greatly appreciated; oftentimes your work can be shortened and made clearer.

\noindent
I also recommend talking with other students, and I am able to give hints to public posts in Piazza.

\noindent
\barsn

\noindent\TAMU{{\sf Hand in at the start of class, Thursday 22 October:}} 


\begin{enumerate}
\setcounter{enumi}{54}


%%%%%%%%%%%%%%%%%%%%%%%%%%%%%%%%%%%%%%%%%%%%%%%%%%%%%%%%%%%%%%%%%%%%%%%%%%%%%%%%%%%%%%%%%%%%%%%%%%%%
\item
  Let $\varphi\colon G \to H$ be a homomorphism of groups and suppose that $K\subset H$ is a normal subgroup of $H$.
  Prove that $\varphi^{-1}(K)$ is a normal subgroup of $G$.
%%%%%%%%%%%%%%%%%%%%%%%%%%%%%%%%%%%%%%%%%%%%%%%%%%%%%%%%%%%%%%%%%%%%%%%%%%%%%%%%%%%%%%%%%%%%%%%%%%%%


%%%%%%%%%%%%%%%%%%%%%%%%%%%%%%%%%%%%%%%%%%%%%%%%%%%%%%%%%%%%%%%%%%%%%%%%%%%%%%%%%%%%%%%%%%%%%%%%%%%%
\item
  Prove or give a counterexample:  If $K$ is a normal subgroup of $G$ and $\varphi\colon G \to H$ is a homomorphism of
  groups,  then the image $\varphi(K)$ of $K$ in $H$ is a normal subgroup of $H$.
%%%%%%%%%%%%%%%%%%%%%%%%%%%%%%%%%%%%%%%%%%%%%%%%%%%%%%%%%%%%%%%%%%%%%%%%%%%%%%%%%%%%%%%%%%%%%%%%%%%%


%%%%%%%%%%%%%%%%%%%%%%%%%%%%%%%%%%%%%%%%%%%%%%%%%%%%%%%%%%%%%%%%%%%%%%%%%%%%%%%%%%%%%%%%%%%%%%%%%%%%
\item
  The symmetric group $S_4$ acts on indeterminates $x_1,x_2,x_3,x_4$ by permuting their indices, and
  this extends to multivariate polynomials $\QQ[x_1,x_2,x_3,x_4]$.
  Specifically
  \[
  \sigma.p(x_1,x_2,x_3,x_4)\ =\ p(x_{\sigma(1)},x_{\sigma(2)},x_{\sigma(3)},x_{\sigma(4)})\,.
  \]
  \begin{enumerate}
  \item
    Find the orbit of each of the following polynomials
    \[
    x_1\;, \qquad x_1+x_2+x_3+x_4\;, \qquad x_1x_2\;, \qquad x_2^2 x_4\,.
    \]
  \item
    Find all fixed points of this action on polynomials of homogeneous degree one.  (That is, polynomials of degree 1 with constant term 0.)
  \item
    Find all fixed points of this action on polynomials of homogeneous degree two.  (That is, polynomials of degree 2 with constant term
    0 and all degree 1 terms also zero.)
  \end{enumerate}
%%%%%%%%%%%%%%%%%%%%%%%%%%%%%%%%%%%%%%%%%%%%%%%%%%%%%%%%%%%%%%%%%%%%%%%%%%%%%%%%%%%%%%%%%%%%%%%%%%%%

%%%%%%%%%%%%%%%%%%%%%%%%%%%%%%%%%%%%%%%%%%%%%%%%%%%%%%%%%%%%%%%%%%%%%%%%%%%%%%%%%%%%%%%%%%%%%%%%%%%%
\item
   Let $B\subset GL(2,\QQ)$ be the set of $2\times 2$ invertible matrices with rational number entries whose lower-left
   entry is zero:
   $\defcolor{B}\vcentcolon=\left\{ \left(\begin{smallmatrix}a&b\\0&c\end{smallmatrix}\right) \mid   a,b,c\in\QQ\ ac\neq0\right\}$,
   and let $N\subset B$ be those matrices with 1s on their diagonal,  
   $\defcolor{N}\vcentcolon=\left\{ \left(\begin{smallmatrix}1&b\\0&1\end{smallmatrix}\right) \mid   b\in\QQ\right\}$.
   \begin{enumerate}
     \item
       Show that both $B$ and $N$ are subgroups of $GL(2,\QQ)$.

     \item
       Prove that $N$ is normal in $B$.

     \item
       Prove that $B/N\simeq \QQ^\times \times \QQ^\times$.

   \end{enumerate}       
%%%%%%%%%%%%%%%%%%%%%%%%%%%%%%%%%%%%%%%%%%%%%%%%%%%%%%%%%%%%%%%%%%%%%%%%%%%%%%%%%%%%%%%%%%%%%%%%%%%%

%%%%%%%%%%%%%%%%%%%%%%%%%%%%%%%%%%%%%%%%%%%%%%%%%%%%%%%%%%%%%%%%%%%%%%%%%%%%%%%%%%%%%%%%%%%%%%%%%%%%
 \item
   Let $G$ be a group and $H$ be a subgroup of $G$.
   Define
   \[
   \defcolor{N_G(H)}\ \vcentcolon=\ \{ g\in G\mid gHg^{-1} = H\}
   \quad\mbox{ and  }\quad
   \defcolor{C_G(H)}\ \vcentcolon=\ \{ g\in G\mid gh=hg\mbox{ for all } h\in H\}\,.
   \]
   These are, respectively, the \demph{normalizer} of $H$ in $G$ and
   the \demph{centralizer} of $H$ in $G$.

   Prove that they are subgroups of $G$.
%%%%%%%%%%%%%%%%%%%%%%%%%%%%%%%%%%%%%%%%%%%%%%%%%%%%%%%%%%%%%%%%%%%%%%%%%%%%%%%%%%%%%%%%%%%%%%%%%%%%

%%%%%%%%%%%%%%%%%%%%%%%%%%%%%%%%%%%%%%%%%%%%%%%%%%%%%%%%%%%%%%%%%%%%%%%%%%%%%%%%%%%%%%%%%%%%%%%%%%%%
\item
  Continuing the previous exercise\ldots
  \begin{enumerate}
  \item
    Prove that $C_G(H)$ is a normal subgroup of $N_G(H)$.
  \item
    Prove that  $N_G(H)/C_G(H)$ is isomorphic to a subgroup of $\Aut(H)$.
  \end{enumerate}
%%%%%%%%%%%%%%%%%%%%%%%%%%%%%%%%%%%%%%%%%%%%%%%%%%%%%%%%%%%%%%%%%%%%%%%%%%%%%%%%%%%%%%%%%%%%%%%%%%%%

%%%%%%%%%%%%%%%%%%%%%%%%%%%%%%%%%%%%%%%%%%%%%%%%%%%%%%%%%%%%%%%%%%%%%%%%%%%%%%%%%%%%%%%%%%%%%%%%%%%%
\item
  Let $G$ be a group and $X$ be a $G$-set.
  Let $g\in G$ and $x\in X$.
  Prove that
  \[
      G_{g.x}\ =\ g G_x g^{-1}\,.
  \]
%%%%%%%%%%%%%%%%%%%%%%%%%%%%%%%%%%%%%%%%%%%%%%%%%%%%%%%%%%%%%%%%%%%%%%%%%%%%%%%%%%%%%%%%%%%%%%%%%%%%


%%%%%%%%%%%%%%%%%%%%%%%%%%%%%%%%%%%%%%%%%%%%%%%%%%%%%%%%%%%%%%%%%%%%%%%%%%%%%%%%%%%%%%%%%%%%%%%%%%%%
\item
  How many distinct octahedral dice are there, up to rotation and reflection?
  
%%%%%%%%%%%%%%%%%%%%%%%%%%%%%%%%%%%%%%%%%%%%%%%%%%%%%%%%%%%%%%%%%%%%%%%%%%%%%%%%%%%%%%%%%%%%%%%%%%%%



\end{enumerate}
%%%%%%%%%%%%%%%%%%%%%%%%%%%%%%%%%%%%%%%%%%%%%%%%%%%%%%%%%%%%%%%%%%%%%%%%%%%%%%%%%%%%%%%%%%%%%%%%%%%%

\end{document}


%%%%%%%%%%%%%%%%%%%%%%%%%%%%%%%%%%%%%%%%%%%%%%%%%%%%%%%%%%%%%%%%%%%%%%%%%%%%%%%%%%%%%%%%%%%%%%%%%%%%
\item

%%%%%%%%%%%%%%%%%%%%%%%%%%%%%%%%%%%%%%%%%%%%%%%%%%%%%%%%%%%%%%%%%%%%%%%%%%%%%%%%%%%%%%%%%%%%%%%%%%%%
